\scanpage{119}\mspage{60}

Du coup il est clair que $A_\lambda$ définit une structure de 2-polyèdre régulier et dont et \deleted{tout}

\[
P_\lambda=\bigl(S,\underset{\substack{\cap\\\unsure{\mathbb P_2|\mathcal P_2|\mathfrak P_2}(S)}}{A_\lambda},
 \underset{\substack{\cap\\\mathbb P_3(S)}}{F_\lambda},\ \text{relations d'incidence évidentes}\bigr)
\tag*{$(43)$}
\]

En fait, on a vu que $G_0^+$ opère par automorphismes de $\Gamma_\lambda$, donc il opère par automorphismes de $P_\lambda$ -- on devine que ce sont des automorphismes directs -- \margin{à y revenir} mais on voudrait que $G_0$ aussi opère \deleted{sur} par $P_\lambda$-automorphismes, pour ceci on va supposer

\[
\begin{aligned}
(T_0\,6)\quad &Z'\text{ commute (non seulement à}\\
&G_0^+,\text{ mais) à tout le groupe }G_0,
\end{aligned}\tag*{$(\deleted{\unsure{42}}\unsure{42|43|44})$}
\]

% slanted
\margin{on pourrait aussi continuer sans cet axiome, en travaillant avec le sous-groupe de $Z'$ commutant à $G_0$ tt entier -- en ignorant \unsure{=|--} l'orientabilité de $P_0$ i.e. laissant tomber l'axiome $T_0\,2$.}

qui implique que $G_0$ opère sur $S$ par $\Gamma_\lambda$-automorphismes, donc par $P_\lambda$-automorphismes. \deleted{On} J'aimerais vérifier que l'hom. canonique injectif

\[
\deleted{\unsure{(44)}}\qquad G_0\lhook\joinrel\longrightarrow\operatorname{Aut}(P_\lambda)\overset{\text{déf}}{=}G_\lambda
\]

est un isom -- comme $P_\lambda\simeq P_0$, on a $G_\lambda\simeq G_0$, et l'isomorphie est évidente si
% archive sheet number: 118
